\section{The rank-deficient averages at the designated primes}

% BEGIN N86V1
\paragraph{N86V1. Good finite series kill both mirabolic character averages.}
Fix a designated residue field $k$, a nonzero vector or covector $a$,
and the character $\theta$ of $k^\times$ prescribed by the member
there. Put $G_r=\mathrm{GL}_r(k)$. Take $H$ to be the subgroup
fixing $\ker a$ pointwise, or the subgroup acting trivially on
$k^r/ka$, respectively. For $\epsilon=1,-1$ put
\[
        m_{H,\epsilon}=\sum_{h\in H}\theta(\det h)^\epsilon h.
\]
Let $\tau$ be a finite irreducible representation in a Green
series whose polynomial omits the rank-one labels of both $\theta$
and $\theta^{-1}$. Then $m_{H,\epsilon}\tau=0$.
Indeed the normalized average is the projection onto vectors
transforming by $\theta^{-\epsilon}\circ\det$ under $H$.
In the hyperplane case, choose coordinates making $\ker a$ the first
$r-1$ axes. The subgroup consists of the unipotent radical of the
$(r-1,1)$ parabolic together with its final $\mathrm{GL}_1$ factor.
In the line case, take the line first and use the $(1,r-1)$ parabolic;
$H$ consists of its radical and its first $\mathrm{GL}_1$ factor.
In either case a nonzero projected vector lies in the radical
invariants and has the specified character on this $\mathrm{GL}_1$.
The complementary $\mathrm{GL}_{r-1}$ commutes with that factor
and acts on its character space. Finite semisimplicity then gives an
irreducible complementary constituent $\sigma$. By finite Frobenius
reciprocity, $\tau$ occurs in the Harish--Chandra induction of
$\sigma$ and $\theta^{-\epsilon}$ in the indicated order.
The cuspidal support of every such constituent contains the
rank-one label of $\theta^{-\epsilon}$, contradicting the
hypothesis. Conjugating coordinates gives every nonzero $a$.

At an original designated place the two character labels are the two
values $\chi_J,\chi_J^{-1}$ in some order, by T86C3. Each tuple
in the good idempotent $e_0$ avoids both values at each such place.
This follows also directly from the good prime-to-$\ell$ root test:
equality of a root with the forbidden value gives equality of their
prime-to-$\ell$ parts for any prime $\ell\ne\operatorname{char}k$.
Use P84C1 to expand $e_0$ as a sum of products of series idempotents.
The average in a single factor annihilates its local series idempotent
by the preceding paragraph, and commutes with all factors in other
places. Thus $m_{H,\epsilon}e_0=0$ in the actual complex group
algebra. In the member's level realization the action of the finite
quotient is by the level-character-corrected translations. Composition
of two such translations is the action of the product: the lift
discrepancy is in the kernel of reduction and its right action cancels
its level-character correction. The recorded finite-group and
coefficient identifications read the distinguished good signature as
this $e_0$ on scalar components. Hence any component fixed by it
is annihilated by all these averages. This proves the source's
$\mathrm{GoodVt}(y)\Rightarrow\mathrm{NoChi}(y)$, the
mirabolic-vanishing input in Y86N2 on the projected space. It does
not require archimedean generation or the regularized integral.
% END N86V1
